\section{Discussion}

The core theorem is the polynomial realization theorem.  Given a squarefree
presentation \(K[T]/(P)\) and an admissible \(a\), it gives an explicit Keller
map with that full Keller fiber.  Monogenicity then yields the corresponding
existence statement for an abstract finite \'etale algebra.  The distinction
matters: the formulas are explicit from \((P,a)\), while the primitive-element
step supplies no canonical presentation.  The classification by rank is a
corollary, but it is not the main point: two finite \'etale algebras of the
same rank can have different field decompositions, Galois actions, signatures,
and local behavior, and the construction preserves all of that data.

Several consequences are immediate.  A squarefree polynomial of degree at
least three becomes an exact inverse polynomial after one translation; an
irreducible polynomial gives a connected full Keller fiber; and a
degree-\(N\) input is realized in \(\A^3\) with coordinate degree at most
\(6N+2\).  Every map
produced by the construction has geometric and arithmetic generic monodromy
\(S_N\).  The primitive action makes every realization map absolutely and
stably compositionally indecomposable.  Thus monodromy supplies a structural
property of the polynomial map, rather than only a generic Galois label.
Over a Hilbertian field, each such fixed map has infinitely many connected
full Keller fibers with splitting-field group \(S_N\).

The cubic seed admits two useful explanations.  Normalized multiplication of
a marked linear and quadratic binary form explains its three marked-root
branches \cite{TaoDigest}; the mixed-sign quotient explains its grading and
contracted locus \cite{Shaska}.  These viewpoints motivate the gauge.  The
load-bearing argument remains the exact identity \(E'=g_1D\) and the boxed
root-to-source reconstruction, which work for every prescribed inverse
polynomial and every commutative test algebra.

Appendix~\ref{app:gauge-identities} retains the coefficientwise expansion
from the compact inverse design to the displayed polynomial coordinates.
Appendix~\ref{app:degree-two} records the scalar-extension and descent
argument for the classical degree-two obstruction, and
Appendix~\ref{app:lean-correspondence} gives the verification table linking
the main statements to their Lean declarations.  Finally,
Appendix~\ref{app:nonproperness} contains the logically separate
nonproperness calculation: the exact reduced nonproperness locus, the
distinguished-target complement, and the complete fiber table on
\(\Pi=0\).

This is a transfer statement, not an inverse Galois theorem.  It does not say
that arbitrary local specifications can always be globalized.  It says that
once a finite \'etale algebra with the desired arithmetic properties exists,
the Keller construction preserves it exactly.  Nor does the monodromy theorem
say that one fixed seed realizes every rank-\(N\) finite \'etale algebra: the
prescribed algebra is used to choose the seed.  Its strengthening is that
every resulting cover, including one built for a highly special algebra, has
maximal generic monodromy.

\paragraph{Relation to earlier work.}
To my knowledge, the following prescribed-algebra statement has not appeared
before: every finite \'etale \(K\)-algebra of rank at least three occurs as a
\emph{full Keller fiber} of a map \(\A^3_K\to\A^3_K\), with explicit
formulas once a squarefree presentation and admissible translation are
supplied.  The rank set \(\{1,3,4,\ldots\}\) follows from this theorem together
with the known degree-two obstruction.  I searched for earlier results up to
24 July 2026, but this priority statement is necessarily qualified: searches
can miss unpublished work or a theorem stated in different language.

The existence of Keller maps in every generic degree is shared terrain.
Gallagher gives explicit all-degree families with complete split rational
fibers \cite{Gallagher}; Migus explicitly uses the available all-degree
constructions as input in the real non-dense-image problem \cite{Migus}.
Shaska studies the grading, quotient geometry, and arithmetic thinness of the
cubic map \cite{Shaska}.  None of these results reverses the quantifiers to
realize an arbitrarily prescribed finite \'etale algebra as the complete
fiber.  Miranda--Neto studies decompositions of fibers already arising from
Keller maps \cite{Miranda}, while Lipton and Markakis connect rational images
with Hilbert irreducibility \cite{LiptonMarkakis}.  The monodromy argument
here uses the classical theorem that a Morse polynomial has symmetric Galois
group over the value-function field \cite[Theorem~4.4.5]{Serre}.

\paragraph{Provenance and AI assistance.}
The recent three-dimensional counterexample was announced by Alp\"oge, who
credited Akhil Mathew for suggesting the problem and Fable for producing the
example \cite{Alpoge}.  Its determinant and collision have since received an
independent Isabelle/HOL formalization \cite{IsabelleAFP}; this paper retains
it only as the cubic seed \(S+S^3\).  Gallagher independently developed an
all-degree atlas \cite{Gallagher}.  Campbell, Razar, and Wright supply the
degree-two obstruction \cite{Campbell,Razar,Wright}.

I used generative AI tools during exploration, symbolic-checker development,
proof review, literature search, Lean formalization, and manuscript editing.
They were useful as fast but fallible collaborators.  I checked the stated
mathematical claims, rewrote the exposition, and take responsibility for the
paper.
